% Source: laplacian.tex
\begin{notechapter}
\chapter{The Laplacian, Local Averages and Graphs}\label{ch:laplacian-average}

\section{Detailed Laplacian-Average Relationship}

\subsection{Step 1: Multivariable Taylor Expansion}
Expand $f(\mathbf{x} + \mathbf{h})$ around $\mathbf{x} = (x,y)$ where 
$\mathbf{h} = (h_x, h_y) = (\epsilon \cos\theta, \epsilon \sin\theta)$:
%
\begin{align*}
f(x + h_x, y + h_y) 
    = f(x,y) 
    + h_x \frac{\partial f}{\partial x}\bigg|_{\mathbf{x}} 
    + h_y \frac{\partial f}{\partial y}\bigg|_{\mathbf{x}}
    + \frac{1}{2}\biggl(
        h_x^2 \frac{\partial^2 f}{\partial x^2}\bigg|_{\mathbf{x}} 
        + 2h_x h_y \frac{\partial^2 f}{\partial x \partial y}\bigg|_{\mathbf{x}}
        + h_y^2 \frac{\partial^2 f}{\partial y^2}\bigg|_{\mathbf{x}} 
    \biggr) 
    + \mathcal{O}(\epsilon^3)
\end{align*}

\subsection{Step 2: Substitute into Circular Average Integral}
Define circular average over radius $\epsilon$ in polar coordinates:
%
\begin{align*}
\text{Avg}(f) 
    &= \frac{1}{2\pi} \int_0^{2\pi} f(x+\epsilon \cos\theta, y+\epsilon \sin\theta) d\theta \\
    &= \frac{1}{2\pi} \int_0^{2\pi} \biggl[
        f(x,y) 
        + \epsilon \cos\theta \frac{\partial f}{\partial x} 
        + \epsilon \sin\theta \frac{\partial f}{\partial y} 
        + \frac{\epsilon^2}{2}\cos^2\theta \frac{\partial^2 f}{\partial x^2}
        + \epsilon^2 \cos\theta \sin\theta \frac{\partial^2 f}{\partial x \partial y} 
        + \frac{\epsilon^2}{2}\sin^2\theta \frac{\partial^2 f}{\partial y^2} 
    \biggr] d\theta 
    + \mathcal{O}(\epsilon^3)
\end{align*}

\subsection{Step 3: Term-by-Term Integration}
Using angular averages:
\begin{align*}
&\int_0^{2\pi} \cos\theta d\theta = 0, \quad 
\int_0^{2\pi} \sin\theta d\theta = 0 \\
&\int_0^{2\pi} \cos^2\theta d\theta = \pi, \quad 
\int_0^{2\pi} \sin^2\theta d\theta = \pi, \\
&\int_0^{2\pi} \cos\theta\sin\theta d\theta = 0
\end{align*}
%
Surviving terms:
\begin{align*}
\text{Avg}(f) &= f(x,y) 
    + \frac{\epsilon^2}{4}\left(\frac{\partial^2 f}{\partial x^2} + \frac{\partial^2 f}{\partial y^2}\right) 
    + \mathcal{O}(\epsilon^3) \\
    &= f(\mathbf{x}) 
    + \frac{\epsilon^2}{4} \nabla^2 f(\mathbf{x}) 
        + \mathcal{O}(\epsilon^3)
\end{align*}
(In fact the remainder is $\mathcal{O}(\epsilon^4)$: the third-order terms involve odd powers of
$\cos\theta$ and $\sin\theta$ and integrate to zero.)

Hence $\nabla^2 f(\mathbf{x}) = \lim_{\epsilon\to 0} \frac{4}{\epsilon^2}\big(\text{Avg}(f) - f(\mathbf{x})\big)$:
the Laplacian measures by how much the average of $f$ on a small circle around $\mathbf{x}$
exceeds the value at the center. It is positive at a local minimum and negative at a local
maximum. (Averaging over the \emph{disk} $B_\epsilon(\mathbf{x})$ instead of the circle gives the
same conclusion with the constant $8$ instead of $4$ in two dimensions.)

\section{Graph Laplacian and Spectral Operations: Detailed Analysis}

\subsection{Graph Fundamentals}
Let $\mathcal{G}=(V,E,W)$ be an undirected weighted graph where:
\begin{itemize}
    \item $V = \{v_1,...,v_n\}$: Set of $n$ nodes
    \item $E \subseteq V \times V$: Set of edges
    \item $W \in \mathbb{R}^{n\times n}$: Weight matrix where $W_{ij} \geq 0$ 
        represents connection strength between nodes $i$ and $j$
\end{itemize}

\subsection{Graph Laplacian Construction}
\begin{enumerate}
    \item \textbf{Degree Matrix} $D$: Diagonal matrix with
        \[
        D_{ii} = \sum_{j=1}^n W_{ij} \quad \text{(total weight connected to node $i$)}
        \]
        
    \item \textbf{Combinatorial Laplacian}:
        \[
        L = D - W
        \]
        This is a symmetric positive semi-definite matrix.

    \item \textbf{Normalized Laplacian} (alternative form):
        \[
        \mathcal{L} = D^{-1/2}LD^{-1/2} = I - D^{-1/2}WD^{-1/2}
        \]
\end{enumerate}

\subsection{Laplacian Action on Node Functions}
For any node function $f: V \to \mathbb{R}$ represented as vector 
$\mathbf{f} = (f_1,...,f_n)^\top \in \mathbb{R}^n$, the Laplacian acts as:
\[
(L\mathbf{f})_i = \sum_{j=1}^n W_{ij}(f_i - f_j)
\]
This measures:
\begin{itemize}
    \item \textbf{Local Variation}: Difference between node $i$'s value 
        and its neighbors
        \item \textbf{Smoothness}: a small value of the quadratic form
        $\mathbf{f}^\top L \mathbf{f} = \frac{1}{2}\sum_{i,j} W_{ij}(f_i - f_j)^2$ indicates that $f$ changes
        slowly over connected nodes
\end{itemize}

\subsection{Spectral Graph Theory Fundamentals}

\begin{theorem}[Spectral Decomposition]
The graph Laplacian admits an orthonormal eigen decomposition:
\[
L = U\Lambda U^\top = \sum_{k=1}^n \lambda_k \mathbf{u}_k \mathbf{u}_k^\top
\]
where:
\begin{itemize}
    \item $\Lambda = \text{diag}(\lambda_1,...,\lambda_n)$: Eigenvalues 
        ($0 = \lambda_1 \leq \lambda_2 \leq ... \leq \lambda_n$)
    \item $U = [\mathbf{u}_1|\cdots|\mathbf{u}_n]$: Orthonormal eigenvectors
    \item $\lambda_k$: Represents "frequency" of eigenmode $\mathbf{u}_k$
\end{itemize}
\end{theorem}

\subsection{Graph Fourier Transform}
Analogous to classical Fourier analysis:
\begin{itemize}
    \item \textbf{Fourier Basis}: Eigenvectors $\{\mathbf{u}_k\}$ form 
        orthonormal basis
    \item \textbf{Fourier Transform}: Decompose signal into spectral components
        \[
        \hat{\mathbf{f}} = U^\top \mathbf{f} \quad \Rightarrow \quad 
        \hat{f}_k = \langle \mathbf{u}_k, \mathbf{f} \rangle
        \]
    \item \textbf{Inverse Transform}: Reconstruct signal from spectrum
        \[
        \mathbf{f} = U\hat{\mathbf{f}} = \sum_{k=1}^n \hat{f}_k \mathbf{u}_k
        \]
\end{itemize}

\subsection{Spectral Graph Convolution: Technical Details}
Given input signal $\mathbf{f}$ and filter $g_\theta$:
\begin{enumerate}
    \item \textbf{Transform to Spectral Domain}:
        \[
        \hat{\mathbf{f}} = U^\top \mathbf{f}
        \]
        
    \item \textbf{Apply Spectral Filter}:
        \[
        \hat{\mathbf{f}}' = g_\theta(\Lambda)\hat{\mathbf{f}} = 
        \text{diag}(g_\theta(\lambda_1),...,g_\theta(\lambda_n))\hat{\mathbf{f}}
        \]
        Common filter types:
        \begin{itemize}
            \item Low-pass: $g_\theta(\lambda) = e^{-\tau\lambda}$ 
                (smooths high-frequency noise)
            \item Band-pass: $g_\theta(\lambda) = \bbone_{[a,b]}(\lambda)$
            \item Learned: $g_\theta(\lambda) = \sum_{k=0}^K \theta_k \lambda^k$
        \end{itemize}
    
    \item \textbf{Inverse Transform}:
        \[
        \mathbf{f}' = U\hat{\mathbf{f}}'
        \]
\end{enumerate}

\subsection{Connection to Continuous Laplacian: Deep Analogy}
\begin{tabular}{p{0.45\textwidth}p{0.45\textwidth}}
\textbf{Continuous Space} & \textbf{Graph Space} \\
\hline
$-\nabla^2 f = -\text{div}(\nabla f)$ & $(L\mathbf{f})_i = \sum_j W_{ij}(f_i-f_j)$ \\
$-\nabla^2 f \propto f - $ local average & $(L\mathbf{f})_i = D_{ii}(f_i - $ neighborhood average$)$ \\
Eigenfunctions of Laplacian form Fourier basis & Laplacian eigenvectors form graph Fourier basis \\
$\hat{f}(\xi) = \int f(x)e^{-i\xi x}dx$ & $\hat{f}_k = \langle \mathbf{u}_k, \mathbf{f} \rangle$ \\
\end{tabular}

\begin{remark}[Sign convention]
With $L = D - W$, the graph Laplacian is the discrete analogue of $-\nabla^2$, not of $\nabla^2$:
$(L\mathbf{f})_i$ is proportional to ``value minus average'', while $\nabla^2 f$ is proportional to
``average minus value''. This is why $L$ is positive semi-definite whereas $\nabla^2$ is
negative semi-definite (e.g.\ $\nabla^2 e^{i\xi x} = -\xi^2 e^{i\xi x}$). For instance, on the path
graph with unit weights, $(L\mathbf{f})_i = 2f_i - f_{i-1} - f_{i+1} \approx -h^2 f''(x_i)$ for samples
$f_i = f(x_i)$ with spacing $h$.
\end{remark}

\subsection{Practical Example: 3-Node Graph}
Consider triangle graph with:
\[
W = \begin{bmatrix}
0 & 1 & 1 \\
1 & 0 & 1 \\
1 & 1 & 0
\end{bmatrix}, \quad
D = \begin{bmatrix}
2 & 0 & 0 \\
0 & 2 & 0 \\
0 & 0 & 2
\end{bmatrix}
\]
Laplacian matrix:
\[
L = D-W = \begin{bmatrix}
2 & -1 & -1 \\
-1 & 2 & -1 \\
-1 & -1 & 2
\end{bmatrix}
\]
Eigenvalues: $\lambda_1=0$, $\lambda_2=3$, $\lambda_3=3$ \\
Eigenvectors: $\mathbf{u}_1=\frac{1}{\sqrt{3}}(1,1,1)^\top$, 
$\mathbf{u}_2=\frac{1}{\sqrt{6}}(2,-1,-1)^\top$, etc.


\section{Complete Graph Laplacian Derivation}

\subsection{Node Functions: Formal Definition}
\begin{definition}[Node Function]
For a graph $\mathcal{G}=(V,E,W)$ with $n=|V|$ nodes, a \textbf{node function} 
is any map $f: V \to \mathbb{R}$ assigning real values to vertices. 
Equivalently, it can be represented as a vector $\mathbf{f} \in \mathbb{R}^n$ 
where $\mathbf{f}_i = f(v_i)$.
\end{definition}

\subsection{Laplacian Action Derivation}
\begin{theorem}[Laplacian Action Formula]
For combinatorial Laplacian $L = D - W$ and node function $\mathbf{f}$:
\[
(L\mathbf{f})_i = \sum_{j=1}^n W_{ij}(f_i - f_j)
\]
\end{theorem}

\begin{proof}
Let $D$ be diagonal with $D_{ii} = \sum_j W_{ij}$. Compute $(D\mathbf{f})_i = 
D_{ii}f_i$ and $(W\mathbf{f})_i = \sum_j W_{ij}f_j$. Then:
\begin{align*}
(L\mathbf{f})_i 
    &= (D\mathbf{f} - W\mathbf{f})_i \\
    &= D_{ii}f_i - \sum_{j=1}^n W_{ij}f_j \\
    &= \left(\sum_{j=1}^n W_{ij}\right)f_i - \sum_{j=1}^n W_{ij}f_j \\
    &= \sum_{j=1}^n W_{ij}(f_i - f_j) \quad \qedhere
\end{align*}
\end{proof}

\subsection{Properties of the Graph Laplacian}
\begin{theorem}[Symmetric Positive Semi-Definiteness]
$L$ is symmetric and $\forall \mathbf{f},\ \mathbf{f}^\top L\mathbf{f} \geq 0$
\end{theorem}

\begin{proof}
Symmetry: $W_{ij}=W_{ji} \Rightarrow (D-W)^\top = D-W$. For quadratic form:
\begin{align*}
\mathbf{f}^\top L\mathbf{f} 
    &= \sum_{i,j} W_{ij}f_i^2 - \sum_{i,j} W_{ij}f_if_j \\
    &= \frac{1}{2}\sum_{i,j} W_{ij}(f_i^2 - 2f_if_j + f_j^2) \\
    &= \frac{1}{2}\sum_{i,j} W_{ij}(f_i - f_j)^2 \geq 0 \quad \qedhere
\end{align*}
\end{proof}

\subsection{Spectral Theorem Application}
\begin{theorem}[Spectral Decomposition]
$\exists$ orthonormal basis $\{\mathbf{u}_k\}_{k=1}^n$ and eigenvalues 
$0 = \lambda_1 \leq \cdots \leq \lambda_n$ such that:
\[
L = \sum_{k=1}^n \lambda_k \mathbf{u}_k \mathbf{u}_k^\top
\]
\end{theorem}

\begin{proof}
Follows from the spectral theorem for symmetric matrices. Since $L$ is real 
symmetric, it can be diagonalized by an orthogonal matrix of eigenvectors with 
real eigenvalues. The quadratic form proof shows all $\lambda_k \geq 0$.
\end{proof}

\subsection{Graph Fourier Transform: Complete Derivation}
\begin{definition}[Graph Fourier Basis]
The eigenbasis $\{\mathbf{u}_k\}$ of $L$ serves as the Fourier basis. The 
Fourier transform $\hat{\mathbf{f}}$ is defined via:
\[
\hat{f}_k := \langle \mathbf{u}_k, \mathbf{f} \rangle = \sum_{i=1}^n u_{ki} f_i
\]
\end{definition}

\begin{theorem}[Convolution Theorem]
For diagonalizable filter $g(L) = U g(\Lambda) U^\top$, convolution satisfies:
\[
g(L)\mathbf{f} = U(g(\Lambda)(U^\top \mathbf{f}))
\]
\end{theorem}

\begin{proof}
Using $L = U\Lambda U^\top$ and $U^\top U = I$:
\begin{align*}
g(L)\mathbf{f} 
    &= U g(\Lambda) U^\top \mathbf{f} \\
    &= U (g(\Lambda) \hat{\mathbf{f}}) \\
    &= \sum_{k=1}^n g(\lambda_k) \hat{f}_k \mathbf{u}_k \quad \qedhere
\end{align*}
\end{proof}

\subsection{Worked Example: 3-Node Cycle Graph}
\begin{example}
Let $\mathcal{G}$ be a triangle graph with:
\[
W = \begin{bmatrix}
0 & 1 & 1 \\
1 & 0 & 1 \\
1 & 1 & 0
\end{bmatrix},\ 
D = \begin{bmatrix}
2 & 0 & 0 \\
0 & 2 & 0 \\
0 & 0 & 2
\end{bmatrix},\ 
L = D-W = \begin{bmatrix}
2 & -1 & -1 \\
-1 & 2 & -1 \\
-1 & -1 & 2
\end{bmatrix}
\]

\textbf{Eigenvalues:} Solve $\det(L-\lambda I)=0$:
\[
\begin{vmatrix}
2-\lambda & -1 & -1 \\
-1 & 2-\lambda & -1 \\
-1 & -1 & 2-\lambda
\end{vmatrix} 
= (3-\lambda)(3-\lambda)(0-\lambda) = 0
\]
Thus $\lambda_1=0$, $\lambda_2=3$, $\lambda_3=3$.

\textbf{Eigenvectors:}
\begin{itemize}
\item $\lambda_1=0$: Solve $L\mathbf{u}=0 \Rightarrow \mathbf{u}_1=\frac{1}{\sqrt{3}}(1,1,1)^\top$
\item $\lambda_2=3$: $(L-3I)\mathbf{u}=0$ means $u_1+u_2+u_3=0$ (a 2-dimensional eigenspace); one choice is $\mathbf{u}_2=\frac{1}{\sqrt{6}}(2,-1,-1)^\top$
\item $\lambda_3=3$: Orthogonal to $\mathbf{u}_1,\mathbf{u}_2$ gives $\mathbf{u}_3=\frac{1}{\sqrt{2}}(0,1,-1)^\top$
\end{itemize}
\end{example}

\subsection{Discrete vs Continuous Laplacian}
\begin{theorem}[Local Average Interpretation]
For small $\epsilon > 0$, the continuous Laplacian satisfies:
\[
\nabla^2 f(\mathbf{x}) = \lim_{\epsilon \to 0} \frac{4}{\epsilon^2}\left[\text{Avg}_{\partial B_\epsilon(\mathbf{x})}(f) - f(\mathbf{x})\right]
\]
while the graph Laplacian gives:
\[
(L\mathbf{f})_i = \sum_{j\sim i} W_{ij}(f_i - f_j) = D_{ii}\left(f_i - \text{Avg}_i(f)\right)
\]
where $\text{Avg}_{\partial B_\epsilon(\mathbf{x})}$ is the average over the circle of radius $\epsilon$ centered at $\mathbf{x}$ (in $\mathbb{R}^2$) and $\text{Avg}_i(f) = \frac{1}{D_{ii}}\sum_j W_{ij}f_j$. Note the opposite signs: $L$ corresponds to $-\nabla^2$.
\end{theorem}

\begin{proof}
For graphs: Rearrange $(L\mathbf{f})_i = D_{ii}f_i - \sum_j W_{ij}f_j = D_{ii}(f_i - \frac{1}{D_{ii}}\sum_j W_{ij}f_j)$. The continuous case follows from the Taylor expansion derivation earlier.
\end{proof}
\end{notechapter}
